\documentclass[11pt,a4paper]{article}

% ── Codificación, fuentes y lenguaje ───────────────────────────────────────────
\usepackage[utf8]{inputenc}
\usepackage[T1]{fontenc}
\usepackage[default,scale=0.95]{sourcesanspro}
\renewcommand{\familydefault}{\sfdefault}
\usepackage[spanish,es-noshorthands,es-tabla]{babel}

% ── Geometría y espaciado ─────────────────────────────────────────────────────
\usepackage[top=2.2cm, bottom=2.5cm, left=2.2cm, right=2.2cm, headheight=15pt]{geometry}
\usepackage{setspace}
\setstretch{1.18}
\usepackage{parskip}

% ── Matemáticas y unidades ────────────────────────────────────────────────────
\usepackage{amsmath}
\usepackage{amssymb}
\usepackage{siunitx}

% ── Circuitos ─────────────────────────────────────────────────────────────────
\usepackage[european, straightvoltages]{circuitikz}

% ── Tablas ────────────────────────────────────────────────────────────────────
\usepackage{booktabs}
\usepackage{tabularx}
\usepackage{array}
\usepackage{longtable}
\usepackage{colortbl}
\usepackage{enumitem}

% ── Colores, cajas e iconos ───────────────────────────────────────────────────
\usepackage[dvipsnames]{xcolor}
\usepackage{tcolorbox}
\tcbuselibrary{skins, breakable}
\usepackage{fontawesome5}
\usepackage{graphicx}

% ── Listados de código (dark-mode infográfico) ────────────────────────────────
%   Estilo base: fondo oscuro con números de línea. Los bloques lstlisting se
%   envuelven automáticamente en una caja tcolorbox redondeada.
\usepackage{listings}

% ── Secciones con barra de color (infografía) ────────────────────────────────
\usepackage{titlesec}
\titleformat{\section}
  {\normalfont\LARGE\bfseries\color{azulNoche}}
  {\colorbox{cyanNeon}{\color{fondoOscuro}\thesection}}{0.7em}{}
  [\vspace{2pt}\color{cyanNeon}\rule{\linewidth}{1.8pt}]
\titlespacing*{\section}{0pt}{24pt}{10pt}
\titleformat{\subsection}
  {\normalfont\large\bfseries\color{azulNoche}}
  {\textcolor{cyanNeon}{\thesubsection}}{0.7em}{}
  [\vspace{1pt}\color{grisLinea}\rule{\linewidth}{0.6pt}]
\titlespacing*{\subsection}{0pt}{14pt}{6pt}
\titleformat{\subsubsection}
  {\normalfont\normalsize\bfseries\color{azulNoche}}
  {\textcolor{cyanNeon}{\thesubsubsection}}{0.7em}{}
\titlespacing*{\subsubsection}{0pt}{10pt}{4pt}

% ── Cabeceras y pies de página ────────────────────────────────────────────────
\usepackage{fancyhdr}
\usepackage{lastpage}
\pagestyle{fancy}
\fancyhf{}
\renewcommand{\headrulewidth}{0pt}
\renewcommand{\footrulewidth}{0pt}
\fancyfoot[C]{%
  \begin{minipage}{\linewidth}
    \vspace{2pt}
    \color{cyanNeon}\rule{\linewidth}{1.2pt}\\[2pt]
    \centering\color{grisTexto}\footnotesize
    Página \thepage\ de \pageref{LastPage}
  \end{minipage}%
}

% ── Hiperenlaces ──────────────────────────────────────────────────────────────
\usepackage[colorlinks=true, linkcolor=azulNoche, urlcolor=cyanNeon]{hyperref}
\sloppy
\emergencystretch 3em

% ── Paleta de colores — tecnológico dark-mode ──────────────────────────────────
\definecolor{fondoOscuro}{HTML}{0D1B2A}     % Dark navy — fondo banda título
\definecolor{verdeTurquesa}{HTML}{004D40}    % Verde turquesa — bandas de marca
\definecolor{azulNoche}{HTML}{1B3A6B}       % Midnight blue — secciones, marcos
\definecolor{azulMedio}{HTML}{2A5298}       % Azul medio — variante de acento
\definecolor{cyanNeon}{HTML}{00C8E0}        % Cyan eléctrico — acentos primarios
\definecolor{verdeSignal}{HTML}{00B887}     % Verde señal — fortalezas, éxito
\definecolor{naranjaVivo}{HTML}{FF7043}     % Naranja vivo — mejoras, advertencias
\definecolor{rojoAlerta}{HTML}{E53935}      % Rojo alerta — errores
\definecolor{grisPapel}{HTML}{F4F6FA}       % Fondo tarjetas claras
\definecolor{grisLinea}{HTML}{C8D0E0}       % Bordes sutiles
\definecolor{grisTexto}{HTML}{6B7A99}       % Texto secundario
\definecolor{amarilloNota}{HTML}{FFB300}    % Notas / advertencias doradas

% Colores específicos para bloques de código (dark-mode)
\definecolor{codigoFondo}{HTML}{1E2D3D}      % Fondo del bloque de código
\definecolor{codigoComentario}{HTML}{7A8FA6} % Comentarios
\definecolor{codigoCadena}{HTML}{FF9D5C}     % Cadenas / strings
\definecolor{codigoKeyword}{HTML}{00C8E0}    % Palabras clave
\definecolor{codigoNumero}{HTML}{00E5A0}     % Números

% ── Banda de título: portada infográfica ──────────────────────────────────────
%   Diseño en 2 capas: banda de color (por defecto fondo oscuro) + franja de
%   acento cyan abajo. Uso: \bandaTitulo[color]{título}{subtítulo}.
%   El icono se puede cambiar con \renewcommand{\iconoBanda}{...} (FontAwesome).
\newcommand{\iconoBanda}{microchip}
\newcommand{\bandaTitulo}[3][fondoOscuro]{%
  \begin{tcolorbox}[
    enhanced,
    colback=#1, colframe=#1,
    boxrule=0pt, arc=8pt, outer arc=8pt,
    width=\linewidth,
    top=18pt, bottom=6pt, left=14pt, right=14pt,
    borderline south={3.5pt}{0pt}{cyanNeon},
  ]
    \begin{minipage}[c]{0.07\linewidth}
      \centering
      \textcolor{cyanNeon}{\fontsize{30}{30}\selectfont\faIcon{\iconoBanda}}
    \end{minipage}%
    \hspace{8pt}%
    \begin{minipage}[c]{0.88\linewidth}
      {\fontsize{20}{24}\selectfont\bfseries\color{white}#2}\par\vspace{5pt}%
      \textcolor{cyanNeon!80!white}{\small\faIcon{angle-right}~#3}%
    \end{minipage}
    \vspace{4pt}
  \end{tcolorbox}%
  \vspace{8pt}%
}

% ── Tarjetas de datos (infografía) ────────────────────────────────────────────
\tcbset{
  tarjetaDato/.style={
    enhanced, breakable,
    arc=6pt, outer arc=6pt,
    colback=grisPapel, colframe=azulNoche,
    boxrule=1pt,
    fonttitle=\bfseries\small, coltitle=white,
    attach boxed title to top left={yshift=-3mm, xshift=6mm},
    boxed title style={
      colback=azulNoche, colframe=azulNoche,
      arc=4pt, boxrule=0pt,
    },
    drop fuzzy shadow=azulNoche!30!white,
    top=8pt, bottom=8pt, left=10pt, right=10pt,
  },
  tarjetaIcono/.style={
    enhanced, breakable,
    arc=6pt, outer arc=6pt,
    colback=white, colframe=grisLinea, boxrule=0.8pt,
    fonttitle=\bfseries\small, coltitle=azulNoche,
    borderline west={3pt}{0pt}{cyanNeon},
    drop fuzzy shadow=azulNoche!25!white,
    top=6pt, bottom=6pt, left=10pt, right=10pt,
  },
  cajaTitulo/.style={
    enhanced, breakable,
    arc=6pt, outer arc=6pt,
    colback=azulNoche!10!grisPapel, colframe=azulNoche,
    boxrule=1pt,
    fonttitle=\bfseries, coltitle=white,
    attach boxed title to top left={yshift=-3mm, xshift=6mm},
    boxed title style={
      colback=azulNoche, colframe=azulNoche,
      arc=4pt, boxrule=0pt,
    },
    drop fuzzy shadow=azulNoche!25!white,
    top=8pt, bottom=8pt, left=10pt, right=10pt,
  },
  cajaContenido/.style={
    enhanced, breakable,
    arc=5pt, outer arc=5pt,
    colback=grisPapel, colframe=grisLinea,
    boxrule=0.8pt,
    fonttitle=\bfseries\small, coltitle=azulNoche,
    top=6pt, bottom=6pt, left=10pt, right=10pt,
  },
  cajaFortaleza/.style={
    enhanced, breakable,
    arc=6pt, outer arc=6pt,
    colback=verdeSignal!8!white, colframe=verdeSignal,
    boxrule=1pt,
    borderline west={4pt}{0pt}{verdeSignal},
    fonttitle=\bfseries\small, coltitle=verdeSignal!70!black,
    drop fuzzy shadow=verdeSignal!20!white,
    top=6pt, bottom=6pt, left=10pt, right=10pt,
  },
  cajaMejora/.style={
    enhanced, breakable,
    arc=6pt, outer arc=6pt,
    colback=naranjaVivo!8!white, colframe=naranjaVivo,
    boxrule=1pt,
    borderline west={4pt}{0pt}{naranjaVivo},
    fonttitle=\bfseries\small, coltitle=naranjaVivo!80!black,
    drop fuzzy shadow=naranjaVivo!20!white,
    top=6pt, bottom=6pt, left=10pt, right=10pt,
  },
  cajaRecomendacion/.style={
    enhanced, breakable,
    arc=6pt, outer arc=6pt,
    colback=cyanNeon!6!white, colframe=cyanNeon!60!azulNoche,
    boxrule=1pt,
    borderline west={4pt}{0pt}{cyanNeon},
    fonttitle=\bfseries\small, coltitle=azulNoche,
    drop fuzzy shadow=azulNoche!20!white,
    top=6pt, bottom=6pt, left=10pt, right=10pt,
  },
}

% ── Ícono + texto (encabezados de sección y elementos) ────────────────────────
\newcommand{\iconotexto}[2]{\textcolor{cyanNeon}{\faIcon{#1}}~#2}

% ── Listados de código: estilo dark + auto-envoltura ─────────────────────────
\lstdefinestyle{estiloCodigo}{
    backgroundcolor=\color{codigoFondo},
    basicstyle=\ttfamily\small\color{white},
    commentstyle=\color{codigoComentario}\itshape,
    stringstyle=\color{codigoCadena},
    keywordstyle=\color{codigoKeyword}\bfseries,
    numberstyle=\tiny\color{codigoComentario},
    numbers=left,
    numbersep=10pt,
    showstringspaces=false,
    breaklines=true,
    breakatwhitespace=false,
    tabsize=4,
    frame=none,
    captionpos=b,
    aboveskip=6pt,
    belowskip=4pt,
    xleftmargin=14pt,
    framexleftmargin=14pt,
    literate=
      {á}{{\'a}}1 {é}{{\'e}}1 {í}{{\'i}}1 {ó}{{\'o}}1 {ú}{{\'u}}1
      {Á}{{\'A}}1 {É}{{\'E}}1 {Í}{{\'I}}1 {Ó}{{\'O}}1 {Ú}{{\'U}}1
      {ñ}{{\~n}}1 {Ñ}{{\~N}}1 {ü}{{\"u}}1 {Ü}{{\"U}}1
      {¿}{{?`}}1 {¡}{{!`}}1
}
\lstdefinestyle{estiloPython}{
    style=estiloCodigo,
    language=Python,
}
\lstset{style=estiloCodigo}
\BeforeBeginEnvironment{lstlisting}{%
  \begin{tcolorbox}[
    enhanced, breakable,
    arc=6pt, outer arc=6pt,
    colback=codigoFondo, colframe=azulNoche,
    boxrule=1pt,
    drop fuzzy shadow=azulNoche!25!white,
    top=2pt, bottom=2pt, left=0pt, right=0pt,
  ]%
}
\AfterEndEnvironment{lstlisting}{\end{tcolorbox}}

% ── Cajas de contenido temáticas ──────────────────────────────────────────────
\newtcolorbox{cajaRecuerda}{
    enhanced, breakable,
    arc=6pt, outer arc=6pt,
    colback=azulNoche!10!grisPapel,
    colframe=azulNoche, boxrule=1pt,
    borderline west={4pt}{0pt}{cyanNeon},
    fonttitle=\bfseries\small\color{white},
    title={\faIcon{info-circle}~Recuerda},
    attach boxed title to top left={yshift=-3mm, xshift=6mm},
    boxed title style={colback=azulNoche, arc=4pt, boxrule=0pt},
    drop fuzzy shadow=azulNoche!25!white,
    left=10pt, right=10pt, top=8pt, bottom=8pt,
}

\newtcolorbox{cajaConcepto}{
    enhanced, breakable,
    arc=6pt, outer arc=6pt,
    colback=verdeSignal!8!white,
    colframe=verdeSignal, boxrule=1pt,
    borderline west={4pt}{0pt}{verdeSignal},
    fonttitle=\bfseries\small\color{white},
    title={\faIcon{lightbulb}~Concepto clave},
    attach boxed title to top left={yshift=-3mm, xshift=6mm},
    boxed title style={colback=verdeSignal!80!black, arc=4pt, boxrule=0pt},
    drop fuzzy shadow=verdeSignal!20!white,
    left=10pt, right=10pt, top=8pt, bottom=8pt,
}

\newtcolorbox{cajaEjemplo}{
    enhanced, breakable,
    arc=6pt, outer arc=6pt,
    colback=cyanNeon!5!white,
    colframe=cyanNeon!70!azulNoche, boxrule=1pt,
    borderline west={4pt}{0pt}{cyanNeon},
    fonttitle=\bfseries\small\color{fondoOscuro},
    title={\faIcon{code}~Ejemplo},
    attach boxed title to top left={yshift=-3mm, xshift=6mm},
    boxed title style={colback=cyanNeon, arc=4pt, boxrule=0pt},
    drop fuzzy shadow=azulNoche!20!white,
    left=10pt, right=10pt, top=8pt, bottom=8pt,
}

\newtcolorbox{cajaEjercicio}{
    enhanced, breakable,
    arc=6pt, outer arc=6pt,
    colback=azulMedio!7!white,
    colframe=azulMedio, boxrule=1pt,
    borderline west={4pt}{0pt}{azulMedio},
    fonttitle=\bfseries\small\color{white},
    title={\faIcon{pencil-alt}~Ejercicio},
    attach boxed title to top left={yshift=-3mm, xshift=6mm},
    boxed title style={colback=azulMedio, arc=4pt, boxrule=0pt},
    drop fuzzy shadow=azulNoche!20!white,
    left=10pt, right=10pt, top=8pt, bottom=8pt,
}

\newtcolorbox{cajaReto}{
    enhanced, breakable,
    arc=6pt, outer arc=6pt,
    colback=naranjaVivo!6!white,
    colframe=naranjaVivo, boxrule=1pt,
    borderline west={4pt}{0pt}{naranjaVivo},
    fonttitle=\bfseries\small\color{white},
    title={\faIcon{fire}~Reto},
    attach boxed title to top left={yshift=-3mm, xshift=6mm},
    boxed title style={colback=naranjaVivo, arc=4pt, boxrule=0pt},
    drop fuzzy shadow=naranjaVivo!20!white,
    left=10pt, right=10pt, top=8pt, bottom=8pt,
}

% ── Indicadores de dificultad ─────────────────────────────────────────────────
\newcommand{\facil}{\textcolor{verdeSignal}{\faIcon{circle}\,\textbf{Fácil}}}
\newcommand{\intermedio}{\textcolor{naranjaVivo}{\faIcon{circle}\,\textbf{Intermedio}}}
\newcommand{\dificil}{\textcolor{rojoAlerta}{\faIcon{circle}\,\textbf{Difícil}}}

% ─────────────────────────────────────────────────────────────────────────────

% ── Cabeceras específicas del documento ──────────────────────────────
\fancyhead[L]{\small\color{azulNoche}\textbf{Servidor MCP S\'ismico} --- Propuesta}
\fancyhead[R]{\small 2026-09-23}

\begin{document}

\renewcommand{\iconoBanda}{satellite}
\bandaTitulo[verdeTurquesa]{Servidor MCP para compartir datos s\'ismicos}{Propuesta para el Prof.\ Rommel Contreras \textbullet\ Colaboraci\'on multi-investigador v\'ia LLMs}

\vspace{6pt}
\begin{tcolorbox}[tarjetaIcono, title={\faIcon{user-graduate}~Resumen para el Prof.\ Rommel}]
Usted puede instalar en su m\'aquina un \textbf{servidor MCP} (Model Context Protocol) que exponga
por red sus cat\'alogo de eventos s\'ismicos y waveforms. Cada investigador, desde su propio
ordenador y con su propia interfaz de IA (opencode, Claude Desktop u otra compatible con MCP),
se conecta remotamente a su servidor y consulta los datos con lenguaje natural. El servidor
\textbf{no requiere APIs de LLM}: solo sirve datos y herramientas; cada colega aporta su modelo.

\medskip
\textbf{Ventaja clave:} ninguna m\'aquina necesitar\'a acceso f\'isico a sus archivos. El dato
vive en un \'unico lugar (su PC) y los modelos de los dem\'as investigadores lo consultan v\'ia IA.
\end{tcolorbox}

% ══ MODELO MENTAL ════════════════════════════════════════════════════
\section{\iconotexto{brain}{C\'omo encaja MCP en la colaboraci\'on}}

MCP es un protocolo \textbf{cliente\,$\leftrightarrow$\,servidor} con dos transportes:

\begin{tcolorbox}[tarjetaDato, title={\faIcon{project-diagram}~Transportes de MCP}]
\begin{tabularx}{\linewidth}{@{}p{3.4cm}XX@{}}
\toprule
\textbf{Transporte} & \textbf{\'Uso} & \textbf{? Servidores remotos} \\
\midrule
\texttt{stdio}        & Servidor en la misma m\'aquina del cliente & No (local) \\
\texttt{Streamable HTTP} & Servidor accesible por red (URL con \texttt{https://}) & S\'i \\
\bottomrule
\end{tabularx}
\end{tcolorbox}

Para colaboraci\'on entre investigadores se usa \textbf{\texttt{Streamable HTTP}}. Punto importante:
el servidor MCP \textbf{no necesita ninguna API de modelo}. Solo expone \emph{datos + herramientas}
(consulta de eventos, waveform, filtros por magnitud/profundidad). Cada investigador corre su propio
cliente que \emph{trae su propio LLM}. Es decir: una m\'aquina en Venezuela aloja los datos; las LLMs
de cada colega se conectan remotamente. El geo-bloqueo de Venezuela aplica a \emph{APIs de LLM}
(se resuelve con OpenRouter o Gemini), no a servir un endpoint HTTP local.

% ══ TOPOLOGÍA ════════════════════════════════════════════════════════
\section{\iconotexto{network-wired}{Topolog\'ia propuesta}}

\begin{center}
\begin{tcolorbox}[cajaContenido, width=0.92\linewidth, title={\faIcon{sitemap}~Red de colaboraci\'on}]
\small
Investigador A \textbullet\ Investigador B \textbullet\ Investigador C\strut\\
\texttt{HTTPS + Bearer token} \,\,$\longrightarrow$\,
\textbf{[\,FastMCP \,---\, Streamable HTTP\,]} \,\,$\longrightarrow$\,
\texttt{/mcp} puerto 8000 $\longrightarrow$ datos s\'ismicos locales\\
\qquad\qquad\qquad\qquad$\uparrow$\\
\qquad\qquad\texttt{cloudflared} (t\'unel) o nginx + Let's Encrypt
\end{tcolorbox}
\end{center}

Cada investigador: \textbf{cliente MCP + su propio modelo}. El Prof.\ Rommel: \textbf{servidor + datos}.

En Venezuela la mayor\'ia de ISPs usan \textbf{CGNAT} (sin IP p\'ublica ni port-forwarding), por lo
que el acceso m\'as sensato y gratuito es \textbf{Cloudflare Tunnel} (\texttt{cloudflared}): otorga
una URL HTTPS sin IP est\'atica. Alternativa cl\'asica: \texttt{nginx} + \textbf{Let's Encrypt} si se
dispone de IP p\'ublica o dominio.

% ══ SERVIDOR ══════════════════════════════════════════════════════════
\section{\iconotexto{server}{Configuraci\'on del servidor (m\'aquina del Prof.\ Rommel)}}

\begin{enumerate}[leftmargin=1.8em, itemsep=4pt]
  \item Instalar el framework:
\begin{lstlisting}
pip install fastmcp
\end{lstlisting}
  \item C\'odigo m\'inimo del servidor:
\begin{lstlisting}[language=Python]
# servidor_sismico.py
from fastmcp import FastMCP

mcp = FastMCP("Sismologia-VE")   # nombre visible para los clientes

@mcp.tool
def eventos(magnitud_min: float = 0.0,
            prof_max_km: float = 300.0) -> list[dict]:
    # Catálogo de eventos sísmicos filtrado por magnitud/profundidad.
    ...  # lee el catálogo local (CSV/Parquet) y filtra

@mcp.tool
def waveform(id_evento: str) -> bytes:
    # Devuelve el waveform del evento id_evento (SAC/mseed).
    ...  # sirve el archivo local

if __name__ == "__main__":
    mcp.run(transport="streamable-http",
            host="0.0.0.0", port=8000)
\end{lstlisting}
  \item Exponer a internet:
\begin{lstlisting}
cloudflared tunnel --url http://127.0.0.1:8000
# devuelve p. ej. https://xxxx.trycloudflare.com
# -> el endpoint MCP es: https://xxxx.trycloudflare.com/mcp
\end{lstlisting}
\end{enumerate}

% ══ CLIENTES ═════════════════════════════════════════════════════════
\section{\iconotexto{laptop}{Configuraci\'on de los clientes (cada investigador)}}

\textbf{opencode} (en \texttt{opencode.json} global o del proyecto):
\begin{lstlisting}[language=json]
{
  "mcp": {
    "sismica": {
      "type": "remote",
      "url": "https://xxxx.trycloudflare.com/mcp",
      "headers": { "Authorization": "Bearer <token-compartido>" },
      "enabled": true
    }
  }
}
\end{lstlisting}

\textbf{Claude Desktop} (\texttt{claude\_desktop\_config.json}):
\begin{lstlisting}[language=json]
{
  "mcpServers": {
    "sismica": {
      "type": "http",
      "url": "https://xxxx.trycloudflare.com/mcp",
      "headers": { "Authorization": "Bearer <token-compartido>" }
    }
  }
}
\end{lstlisting}

El token se reparte entre los investigadores por canal cifrado; para un grupo cerrado es suficiente
y puede rotarse peri\'odicamente.

% ══ CONSIDERACIONES ══════════════════════════════════════════════════
\section{\iconotexto{shield-alt}{Consideraciones}}

\begin{tcolorbox}[cajaRecomendacion, title={\faIcon{lock}~Seguridad}]
Exigir \textbf{HTTPS} (el t\'unel lo da gratis) y \textbf{Bearer token}; no exponer el puerto 8000
crudo a internet. Para un grupo como una comunidad de investigadores es una soluci\'on pr\'actica.
\end{tcolorbox}

\begin{tcolorbox}[cajaRecomendacion, title={\faIcon{memory}~Concurrencia y carga}]
Varios clientes golpean el servidor a la vez. Si los waveforms son pesados, cachear en disco o servir
una versi\'on decimada por defecto; a menor volumen de datos transferidos, menor consumo de tokens de
los modelos y menor uso de su RAM (recordar que la PC del receptor tiene 4 GB).
\end{tcolorbox}

\begin{tcolorbox}[cajaRecomendacion, title={\faIcon{cloud-upload-alt}~Disponibilidad y respaldo}]
Si la m\'aquina del Prof.\ Rommel se apaga, los datos dejan de estar disponibles. Opci\'on robusta:
publicar el cat\'alogo en GitHub o Hugging Face Dataset, de modo que cada investigador apunte su MCP a
su propio clon y el servidor central solo orqueste los metadatos.
\end{tcolorbox}

\begin{tcolorbox}[cajaMejora, title={\faIcon{exclamation-triangle}~Verificaci\'on del paquete}]
Igual que ocurri\'o con \texttt{@platformio/mcp} (que no exist\'ia; el real era \texttt{platformio-mcp}),
conviene confirmar el nombre y versi\'on del paquete MCP en el \'indice antes de asumir:
\texttt{pip index versions mcp} y revisar PyPI.
\end{tcolorbox}

% ══ CIERRE ═══════════════════════════════════════════════════════════
\vspace{10pt}
\begin{tcolorbox}[cajaConcepto]
\textbf{Pr\'oximo paso acordado:} estudiar la construcci\'on de un \textbf{prototipo funcional en el
repositorio ELECTRONICA} siguiendo la arquitectura de 3 capas --- directiva
\texttt{directives/servidor\_sismico.yaml} + orquestador + \texttt{execution/servidor\_sismico.py}
con un cat\'alogo sint\'etico de ejemplo --- para que el Prof.\ Rommel lo clone y solo cambie la ruta a
sus datos reales.
\end{tcolorbox}

\vspace{12pt}
\begin{center}
\footnotesize\color{grisTexto}
Documento generado desde \texttt{ELECTRONICA} (\textbullet\ estilo infogr\'afico propio). \textbullet\ 2026-09-23.
\end{center}

\end{document}
